% si_maps.tex : generated by tools/make_si_maps.py (do not edit by hand).
% Supplementary Figs S1-S6 (former S10-S15): 1 km ALT maps (display resolution, no accuracy claim) and the XK/XL figures.
% S1 (former S10): Alaska_ALT_map_v3.pdf (file time 2026-10-05 13:03); caption from Alaska_ALT_map_v3_legend.md
\begin{figure}[p]\centering
\IfFileExists{../../../outputs/figures/paper/v3_restructure/maps/Alaska_ALT_map_v3.pdf}{\includegraphics[width=\linewidth,height=0.72\textheight,keepaspectratio]{../../../outputs/figures/paper/v3_restructure/maps/Alaska_ALT_map_v3.pdf}}{\fbox{\parbox[c][45mm][c]{0.85\linewidth}{\centering Supplementary Fig.~S1 in preparation}}}
\caption{\textbf{Active-layer thickness (ALT) of the Alaskan permafrost zone at 1 km display resolution: recalibrated Stefan model, residual machine-learning correction and final map.} a, Recalibrated Stefan ALT, E·√TDD with E = 1.617 (least squares on all 13,606 label cells, shrunk towards the source coefficient). b, Residual ML correction λ·g (final minus recalibrated Stefan; positive, deeper), where g is a CatBoost residual model on 25 covariates (mean of two seeds) and λ = 0.5 was selected from {0.25, 0.5, 1.0} by 0.5° block five-fold cross-validation. c, Final ALT, a plus b. a and c share one colour scale (35–75 cm, 1st–99th percentile of both maps). d, Number of covariates outside the 0.5–99.5 percentile range of the training labels (hatching density), with label cells as dots; 73\% of mapped cells have at least one, mostly terrain roughness, slope and topographic position. The 90\% prediction interval of c is ±20.2 cm everywhere (constant-width split conformal on block cross-validation residuals; cross-validated RMSE 13.8 cm, recalibrated Stefan 14.3 cm; lower bound truncated at 0 cm in 5,384 cells). Its nominal coverage is an average over cells similar to the labels and is not guaranteed in hatched cells, so the width is not mapped. The 1 km grid is a display resolution: air thawing index and climate covariates come from ERA5-Land (0.1°, about 4–6 × 11 km here), soil covariates from SoilGrids windows at about 5 km, and the residual correction acts mainly between climate-grid cells (main text, Fig.~6c). Cells are drawn as their native 0.009° rectangles. Domain: land (Natural Earth 10 m) with ESA CCI permafrost fraction of at least 50\% (1997–2021 mean), 59–71.5° N, 168–141° W; 858,517 of 899,633 cells mapped. White, sea or water (Global Surface Water occurrence of at least 50\%); grey, missing CCI ALT or soil thaw index. NAD83 Alaska Albers (EPSG:3338); inset, location. Data: ABoVE ALT v2 (Moore et al. 2025, ORNL DAAC) and GTN-P/CALM (Streletskiy et al. 2025) labels; ERA5-Land (Copernicus C3S); SoilGrids 2.0 (ISRIC); Copernicus DEM GLO-30 (ESA, AWS Open Data); ESA CCI Permafrost v4.0; JRC Global Surface Water; Natural Earth; Cartopy 0.23.}\label{fig:S10}
\end{figure}
\clearpage
% S2 (former S11): Lena_ALT_map_v3.pdf (file time 2026-10-05 13:03); caption from Lena_ALT_map_v3_legend.md
\begin{figure}[p]\centering
\IfFileExists{../../../outputs/figures/paper/v3_restructure/maps/Lena_ALT_map_v3.pdf}{\includegraphics[width=\linewidth,height=0.72\textheight,keepaspectratio]{../../../outputs/figures/paper/v3_restructure/maps/Lena_ALT_map_v3.pdf}}{\fbox{\parbox[c][45mm][c]{0.85\linewidth}{\centering Supplementary Fig.~S2 in preparation}}}
\caption{\textbf{Active-layer thickness (ALT) of the Lena River Delta at 1 km display resolution: recalibrated Stefan model, residual machine-learning correction and final map.} a, Recalibrated Stefan ALT, E·√TDD with E = 1.435 (least squares on all 3,037 Lena label cells, shrunk towards the source coefficient 1.619). b, Residual ML correction 0.25·g (final minus recalibrated Stefan; positive, deeper), where g is a CatBoost residual model (mean of two seeds) fitted to residuals of source-region cells beyond a 100 km buffer and of the Lena labels. c, Final ALT, a plus b. a and c share one colour scale (35–50 cm, 1st–99th percentile of both maps). d, Number of covariates outside the 0.5–99.5 percentile range of the training rows of c (source cells and Lena labels), with label cells as dots; 84\% of mapped cells have at least one, mostly ERA5-Land coldest-month temperature, mean annual air temperature and freezing index. Relative to source cells alone, all have three or more. The 90\% prediction interval of c is ±25.4 cm everywhere (constant-width split conformal on 0.5° block five-fold cross-validation residuals within the Lena labels; cross-validated RMSE 20.8 cm, recalibrated Stefan 21.5 cm). Its nominal coverage is an average over cells similar to the labels and is not guaranteed in hatched cells, so the width is not mapped. The 1 km grid is a display resolution: air thawing index and climate covariates come from ERA5-Land (0.1°, about 3 × 11 km here), soil covariates from SoilGrids windows at about 5 km, and the residual correction acts mainly at the climate-grid scale (main text, Fig.~6c); the horizontal bands in a follow ERA5-Land rows. Cells are drawn as their native 0.009° rectangles. Grid: 71.5–73.6° N, 123.3–130.1° E; 38,086 mapped cells. White, sea or water (Global Surface Water occurrence of at least 50\%); grey, missing CCI ALT. Polar stereographic projection (central meridian 126.7° E, true scale at 70° N); inset, location. Data: ALLena thaw-depth compilation (Veremeeva et al. 2025, PANGAEA); ERA5-Land (Copernicus C3S); SoilGrids 2.0 (ISRIC); Copernicus DEM GLO-30 (ESA, AWS Open Data); ESA CCI Permafrost v4.0; JRC Global Surface Water; Natural Earth 10 m; Cartopy 0.23.}\label{fig:S11}
\end{figure}
\clearpage
% S3 (former S12): XL_display_resolution.pdf (file time 2026-10-05 13:04); caption from XL_display_resolution_legend.md
\begin{figure}[p]\centering
\IfFileExists{../../../outputs/figures/paper/v3_restructure/maps/XL_display_resolution.pdf}{\includegraphics[width=\linewidth,height=0.72\textheight,keepaspectratio]{../../../outputs/figures/paper/v3_restructure/maps/XL_display_resolution.pdf}}{\fbox{\parbox[c][45mm][c]{0.85\linewidth}{\centering Supplementary Fig.~S3 in preparation}}}
\caption{\textbf{The same residual ML ALT map at its 1 km cells and as ERA5-Land 0.1° cell means.} a, b, Alaska, 50 × 50 km around the 0.5° block with the most label cells (69.1° N, 148.8° W; 2,892 cells); white lines, ERA5-Land 0.1° cell boundaries. c, d, Lena Delta, full grid. a and c show the 1 km cells; b and d colour each cell with the mean of the mapped cells in its ERA5-Land 0.1° cell (boundaries at grid point ± 0.05°). All cells are drawn as their native 0.009° rectangles without resampling. One colour scale per region from the 1st–99th percentile of both panels (Alaska 42–54 cm, Lena Delta 38–48 cm), narrower than in Supplementary Figs~S1 and S2 so that within-cell detail is visible. The 0.1° cell means explain 97\% (Alaska, whole domain) and 61\% (Lena Delta) of the variance of the 1 km map; the standard deviation of the 1 km values around their 0.1° cell mean is 1.8 cm in Alaska (total 9.5 cm) and 1.1 cm in the Lena Delta (total 1.8 cm). The 1 km map is a display resolution; its information resolution is bound to the climate input (ERA5-Land, about 10 km) and the soil input (SoilGrids, about 5 km). Within-cell detail comes from the terrain (Copernicus DEM 30 m windows) and CCI ALT inputs. Projections as in Supplementary Figs~S1 and S2 (NAD83 Alaska Albers; polar stereographic, central meridian 126.7° E). Data as in Supplementary Figs~S1 and S2; Cartopy 0.23.}\label{fig:S12}
\end{figure}
\clearpage
% S4 (former S13): XL_Alaska_product_differences.pdf (file time 2026-10-05 13:03); caption from XL_Alaska_product_differences_legend.md
\begin{figure}[p]\centering
\IfFileExists{../../../outputs/figures/paper/v3_restructure/maps/XL_Alaska_product_differences.pdf}{\includegraphics[width=\linewidth,height=0.72\textheight,keepaspectratio]{../../../outputs/figures/paper/v3_restructure/maps/XL_Alaska_product_differences.pdf}}{\fbox{\parbox[c][45mm][c]{0.85\linewidth}{\centering Supplementary Fig.~S4 in preparation}}}
\caption{\textbf{The 1 km Alaska ALT map compared with four ALT products and with the recalibrated Stefan model.} a–e, ALT of the residual ML map (a, as in Supplementary Fig.~S1) and of CCI v5 (1997–2023 mean, b), Wei 2026 (2000–2024 mean, c), Aalto 2018 (2000–2014 baseline, d) and Yi-Kimball (2001–2015 mean, e) on one colour scale (0–240 cm, 2nd–98th percentile of all five maps over mapped cells). Grey text, area mean over mapped cells with a product value and Pearson correlation r with the residual ML map. f, Recalibrated Stefan − residual ML on its own scale (±10 cm, the scale of panel b of Supplementary Figs~S1 and S2 with the sign reversed). g–j, Product − residual ML (positive, product deeper) on one scale (±190 cm, 98th percentile of the absolute differences of g–j). Products are read at the 1 km cell centres (nearest pixel, period mean). Mean differences are +49.4 cm (CCI v5), +14.7 (Wei 2026), +22.6 (Aalto 2018), +49.3 (Yi-Kimball) and +1.4 cm (Stefan). Large differences with CCI v5 and Yi-Kimball are widespread in Interior Alaska; the largest differences with Wei 2026 lie at about 61–63° N, 141–150° W; Aalto 2018 is shallower than the map at about 62–64° N west of 148° W. These maps carry no accuracy statement; accuracy is compared only at label cells (Supplementary Table S2 and Supplementary Fig.~S6). The 1 km map is a display resolution; its information resolution is bound to the climate input (ERA5-Land, about 10 km) and the soil input (SoilGrids, about 5 km). Grey cells, no product value (Yi-Kimball covers 86\% of mapped cells). NAD83 Alaska Albers (EPSG:3338). Data: CCI v5 (ESA CCI Permafrost, CEDA, doi:10.5285/a6fbedd8ee5b472c8e84e55f746c1704); Wei et al. 2026 (Zenodo 21667583, CC BY 4.0); Aalto et al. 2018 (Zenodo 5003804, CC0); Yi and Kimball (ORNL DAAC 1760); covariates and labels as in Supplementary Fig.~S1; Natural Earth 10 m; Cartopy 0.23.}\label{fig:S13}
\end{figure}
\clearpage
% S5 (former S14): XL_Lena_product_differences.pdf (file time 2026-10-05 13:04); caption from XL_Lena_product_differences_legend.md
\begin{figure}[p]\centering
\IfFileExists{../../../outputs/figures/paper/v3_restructure/maps/XL_Lena_product_differences.pdf}{\includegraphics[width=\linewidth,height=0.72\textheight,keepaspectratio]{../../../outputs/figures/paper/v3_restructure/maps/XL_Lena_product_differences.pdf}}{\fbox{\parbox[c][45mm][c]{0.85\linewidth}{\centering Supplementary Fig.~S5 in preparation}}}
\caption{\textbf{The 1 km Lena Delta ALT map compared with three ALT products and with the recalibrated Stefan model.} a–d, ALT of the residual ML map (a, as in Supplementary Fig.~S2) and of CCI v5 (1997–2023 mean, b), Wei 2026 (2000–2024 mean, c) and Aalto 2018 (2000–2014 baseline, d) on one colour scale (25–75 cm, 2nd–98th percentile of all four maps over mapped cells). Grey text, area mean over mapped cells with a product value and Pearson correlation r with the residual ML map. e, Recalibrated Stefan − residual ML on its own scale (±10 cm, the scale of panel b of Supplementary Figs~S1 and S2 with the sign reversed). f–h, Product − residual ML (positive, product deeper) on one scale (±35 cm, 98th percentile of the absolute differences of f–h). Products are read at the 1 km cell centres (nearest pixel, period mean). Mean differences are −0.9 cm (CCI v5), +16.5 (Wei 2026, 94\% of cells with values), +13.0 (Aalto 2018) and −1.1 cm (Stefan). CCI v5 is shallower than the map in the northern and eastern delta and deeper in the south-western part of the grid; Wei 2026 and Aalto 2018 are deeper over most of the grid. These maps carry no accuracy statement; accuracy is compared only at label cells. The 1 km map is a display resolution; its information resolution is bound to the climate input (ERA5-Land, about 10 km) and the soil input (SoilGrids, about 5 km). Grey cells, no product value. Polar stereographic projection (central meridian 126.7° E, true scale at 70° N). Data: CCI v5 (ESA CCI Permafrost, CEDA); Wei et al. 2026 (Zenodo 21667583, CC BY 4.0); Aalto et al. 2018 (Zenodo 5003804, CC0); covariates and labels as in Supplementary Fig.~S2; Natural Earth 10 m; Cartopy 0.23.}\label{fig:S14}
\end{figure}
\clearpage
% S6 (former S15): XK_support_scale.pdf (file time 2026-10-05 13:04); caption from XK_support_scale_legend.md
\begin{figure}[p]\centering
\IfFileExists{../../../outputs/figures/paper/v3_restructure/maps/XK_support_scale.pdf}{\includegraphics[width=\linewidth,height=0.72\textheight,keepaspectratio]{../../../outputs/figures/paper/v3_restructure/maps/XK_support_scale.pdf}}{\fbox{\parbox[c][45mm][c]{0.85\linewidth}{\centering Supplementary Fig.~S6 in preparation}}}
\caption{\textbf{Prediction error of the recalibrated Stefan model, the residual ML model and four ALT products after aggregation to larger grid cells.} Model values are within-region out-of-block predictions from 0.5° block five-fold cross-validation; the Stefan coefficient is refitted on the training folds, and λ of the residual model is selected by an inner block cross-validation within the training folds. Product values are period means at the label cells: CCI v4 (1997–2021), CCI v5 (1997–2023), Wei 2026 (2000–2024) and, for Alaska, Yi-Kimball (2001–2015). Labels and predictions are averaged in grid cells of 0.05°, 0.1° and 0.25° whose boundaries follow the ERA5-Land cells (0.1° = ERA5-Land cell); grid cells with fewer than three label cells are dropped; 1 km denotes single label cells. Grey numbers, grid cells per size. a–c, Cell-weighted RMSE (logarithmic axis) on the cells where all shown values exist (Alaska 10,605 cells in 60 blocks, Lena Delta 1,342 in 15, Canada 745 in 34); bands, 95\% block bootstrap intervals for the two models (1,000 resamples of 0.5° blocks). d–f, RMSE of the recalibrated Stefan model minus RMSE of the residual ML model on all label cells with model values (Alaska 13,606, Lena Delta 2,958, Canada 747), cell-weighted (filled) and block-equal (open) with 95\% block bootstrap intervals. In Alaska the cell-weighted difference is 0.50 cm (0.18 to 1.10) for single cells and 1.54 to 1.67 cm (lower bounds 0.50 to 0.60) at 0.05°, 0.1° and 0.25°. In Canada it is 2.62 cm (0.00 to 4.42) for single cells and 0.96 to 1.35 cm at larger sizes, with intervals that include zero. In the Lena Delta all intervals include zero. Removing cells within 5 km of Wei 2026 training sites changes the Alaska differences by at most 0.17 cm (data/processed/xbatch/XK\_support\_scale/). Products are compared only as RMSE differences at the same grid size.}\label{fig:S15}
\end{figure}
\clearpage
